\chapter{Conclusion}
\label{ch:conclusion}

This project undertook a systematic, comparative evaluation of three conformal prediction frameworks,
Inductive Conformal Prediction, Conformal Risk Control, and Learn then Test,
within the safety-critical context of multi-label chest X-ray classification.
The investigation sought to move beyond standard coverage metrics to control for specific,
clinically meaningful diagnostic errors,
providing a clear analysis of the trade-offs between statistical guarantees and predictive utility.

The findings confirm that while baseline class-conditional ICP provides valid per-label coverage,
it is unsuitable for risk-sensitive applications.
It yielded a high False Negative Rate (FNR) of approximately 43\%,
demonstrating that a simple coverage guarantee does not protect
against the critical error of missing pathologies.
In contrast, CRC proved highly effective at controlling the FNR by
consistently keeping the average risk at or below the desired 10\% target,
demonstrating its value in creating safer clinical AI systems.

A key finding emerged from the attempt to implement a novel,
clinically-weighted loss function within CRC.
The experiment revealed a critical limitation of the framework;
instead of selectively reducing risk for high-priority pathologies,
the weighted loss increased the FNR across all classes.
This result occurred because the single, global calibration parameter became overly conservative,
degrading overall performance.
This highlights that modifying the loss function in standard CRC
can have unexpected and system-wide consequences.

The LTT framework successfully addressed the primary limitations of CRC.
In controlling the FNR, LTT provided a robust high-probability guarantee,
ensuring that the risk ceiling was violated in fewer than 10\% of trials.
This is a far stronger assurance than CRC's guarantee in expectation
and is better suited for clinical use,
although it came at the cost of slightly larger, less precise prediction sets.
Crucially, LTT demonstrated its unique capability to control non-monotone risks
by successfully managing the False Discovery Rate (FDR) at various target levels.
This extends its utility to addressing important clinical issues like alarm fatigue,
a task that is fundamentally beyond the scope of CRC.
The robustness of these frameworks was confirmed across multiple datasets and model architectures,
and they also proved effective as diagnostic tools for identifying model-task mismatch.

Building on the findings presented, several promising avenues for future research are apparent.

\begin{itemize}
    \item \textbf{Group-Conditional Risk Control}:
    As mentioned in the report, a critical next step is to address potential algorithmic bias
    by extending these frameworks to offer group-conditional guarantees.
    This would involve controlling risks like the FNR separately
    for different demographic subgroups,
    ensuring equitable performance and safety for all patient populations.

    \item \textbf{Multi-Dimensional Risk Control with LTT}:
    The failure of the weighted loss function in CRC highlighted the limitation
    of a single calibration parameter.
    Future work can leverage the LTT framework with a multi-dimensional lambda parameter.
    This would enable granular risk control by assigning different risk targets
    to different pathologies based on their clinical severity,
    allowing the system to be simultaneously strict
    on critical findings and more lenient on minor ones.

    \item \textbf{Simultaneous Control of Multiple Risks}:
    Clinical decision-making often requires balancing competing objectives.
    An important extension would be to configure the LTT framework
    to simultaneously control both the FNR and the FDR.
    By defining a multi-risk objective,
    the system could be calibrated to find prediction sets that,
    with high probability, both minimize missed pathologies
    and avoid overwhelming clinicians with excessive false alarms,
    further enhancing its practical utility and trustworthiness.
\end{itemize}

In conclusion, this work illuminates a clear hierarchy of conformal methods.
CRC offers a significant improvement for controlling specific monotone risks,
but its guarantee is limited to the average case
and its mechanism is sensitive to global parameter changes.
LTT provides the most powerful and flexible approach,
offering stronger high-probability guarantees for any risk function.
The choice between them presents a clear trade-off between
the strength of the statistical guarantee and the precision of the clinical output.
